\section*{\centering \fontsize{35}{45}\selectfont Creating a Knight}
\vspace*{0.4 cm}

\begin{multicols*}{2}

\subsection*{Virtues}

Each knight has three attributes that describe his body, mind, and soul.

\textbf{Vigor.}
Strong limbs, firm hands, powerful lungs.

\textbf{Clarity.}
Keen instinct, lucid mind, shrewd eyes.

\textbf{Spirit.}
Charming tongue, iron will, fierce heart.

Knight's creation starts by rolling \texttt{d6 + d12} four times; record not only the sum, but also individual values of each die.
Each of those rolls determines one of the knight's \texttt{Virtues}; you choose which is which and discard the remaining one. You can use the character sheet on page \pageref*{sec:ch_sheet}.

Then, determine the knight's capacity for \texttt{Injury}(\texttt{Injury Slots}), \texttt{Focus} and \texttt{Stress} according to the steps below.

\colbox{
\begin{itemize}[leftmargin=*]
    \item[\ding{170}] Half, rounded up, of \texttt{d12} you assigned to \texttt{Vigor} is the amount of \texttt{Severe Injuries} the knight can survive before dying.
    \item[\textdagger] Half, rounded up, of \texttt{d6} you assigned to \texttt{Vigor} is the amount of \texttt{Lethal Injuries} the knight can survive before dying.
    \item[\ding{84}] The value of \texttt{d6} you assigned to \texttt{Clarity} is maximum \texttt{Focus} the knight has.
    \item[\ding{118}] Half, rounded up, of \texttt{d12} you assigned to \texttt{Spirit} is maximum \texttt{Stress} the knight can endure before suffering a \texttt{Trauma}, minimum is \texttt{2}.
\end{itemize}
}
\paragraph{Guard.} Roll \texttt{d6} to determine the maximum \texttt{Guard(GD)} the knight has.

\paragraph{Choose a Passion. } Something dear to the knight's heart, or something weird that allows them to shrug off a failure.
Indulging a Passion restores \texttt{Spirit} and \texttt{Spellweaving} dice (p.\pageref*{sec:spellweaving}). Consult with your Referee if in doubt.

\newcolumn

\paragraph{Saves.}

To pass a Save, roll a \texttt{d20} equal to or below the
relevant \texttt{Virtue}. Failure means negative
consequences, not always a failed action.

\paragraph{Feats. } Knight knows all three \texttt{Feats}: \texttt{Smite}, \texttt{Focus} and \texttt{Deny}. 
They can perform them during combat, details of which are on page \pageref*{sec:combat}.

\subsection*{Injury Slots}

When a knight suffers an \texttt{Injury}, fill an \texttt{Injury slot} following the steps described below.

\colbox{
\begin{enumerate}[nosep, leftmargin=*]
    \item Determine if injury is \texttt{Severe} or \texttt{Lethal}
    \item Briefly describe an injury 
    \item Record the damage die of the attack which caused the injury
    \item Roll \texttt{Affliction}(see p.\pageref*{sec:affliction}) and record it
\end{enumerate}
}
Each injury could be:
\begin{itemize}[nosep, leftmargin=*]
    \item \textbf{Untreated} \texttt{injury} applies \texttt{Affliction}(p.\pageref*{sec:injuries})
    \item \textbf{Treated}(p.\pageref*{sec:combat}, \pageref*{sec:healing}) \texttt{injury} occupies the slot, but does not apply \texttt{Affliction}
    \item \textbf{Healed} \texttt{injury} is removed entirely, but leaves a Quirk(p.\pageref*{sec:quirks}) behind.
\end{itemize}
If a knight suffers an injury and all corresponding slots are full, fill one with greater lethality. If they are out of those, the knight is \textbf{Slain}.

\subsection*{Spellcasting}

The Realm is not strange to magic and some knights employ it during their endeavours. 

A knight might be a \textbf{Spellweaver}(see p.\pageref*{sec:spellweaving}) or a \textbf{Witch}(see p.\pageref*{sec:witchcraft}). 
If you want to tap into those powers choose a suitable trait, see page\pageref*{sec:traits}.

\subsection*{Weapons \& Gear}

When picking gear for a knight, consult with the Referee, but the rule of thumb is as follows: you have a budget of 12 points; these points can be spent on weapons and armour. 

Weapon costs as much as its damage dice. For example, shortsword(2d6 damage) costs 12 points, knife(d6 damage) - 6 points. Each piece of gear providing armour costs 2 per armour point, e.g. a shield costs 6 points. 
Usually it makes sense to limit starting equipment(p.\pageref*{sec:gear}) to be common rarity with a single uncommon item as a prized possession of your knight.

\end{multicols*}
